\documentclass{article}

\usepackage{tabularx}
\usepackage{booktabs}

\title{Reflection and Traceability Report on \progname}

\author{\authname}

\date{}

\input{../Comments}
%% Common Parts

\newcommand{\progname}{Industrial Plant Designer} % PUT YOUR PROGRAM NAME HERE
\newcommand{\authname}{Team 18, Chemware Engineering
\\ Kishor Pandya
\\ Dhrumil Pandya
\\ Madhu Sivapragasam
\\ Rasik Pokharel} % AUTHOR NAMES                  

\usepackage{hyperref}
    \hypersetup{colorlinks=true, linkcolor=blue, citecolor=blue, filecolor=blue,
                urlcolor=blue, unicode=false}
    \urlstyle{same}
                                


\begin{document}

\maketitle


\section{Changes in Response to Feedback}

\subsection{SRS and Hazard Analysis/POC}

    
    \subsubsection*{Feedback Label: Add filters for node variables}
    \textbf{Source:} Prof Mahalec after POC demo \\
    \textbf{Issue Link:} \href{https://github.com/MadhuranS/capstone/issues/97}{link} \\
    \textbf{PR Link:} \href{https://github.com/MadhuranS/capstone/pull/130}{link} \\
    \textbf{Summary:} Professor Mahalec indicated that he wanted a filter for nodes variables due to the exisiting list of variables overcrowding the UI. In order to simplify the UI and improve usability and readability a filter was needed. 

    \subsubsection*{Feedback Label: Add linter and formatting setup for repository}
    \textbf{Source:} Prof Smith feedback after POC \\
    \textbf{Issue Link:} \href{https://github.com/MadhuranS/capstone/issues/161}{link} \\
    \textbf{PR Link:} \href{https://github.com/MadhuranS/capstone/pull/132}{link} \\
    \textbf{Summary:} Professor Smith indicated that we should have some sort of linting and code formatting tools setup in our repository after the POC demo. To address this we decided to use prettier and ESLint. 

    \subsubsection*{Feedback Label: Add completion check for all diagrams}
    \textbf{Source:} Prof Mahalec's grad student after POC demo \\
    \textbf{Issue Link:} \href{https://github.com/MadhuranS/capstone/issues/166}{link} \\
    \textbf{PR Link:} \href{https://github.com/MadhuranS/capstone/pull/169}{link} \\
    \textbf{Summary:} Professor Mahalec's grad student, Farbod, indicated that there should be some sort of method in place to check if a diagram is a valid diagram or not (valid as free of errors). To do This is because a diagram that does not indicate potential problems lacks robustness and correctness and so it's important to add visibility into the validity. To do so we added a completion check to the diagrams. 


    \subsubsection*{Feedback Label: Add ability to select stream contents}
    \textbf{Source:} Professor Mahalec after POC \\
    \textbf{Issue Link:} \href{https://github.com/MadhuranS/capstone/issues/98}{link} \\
    \textbf{PR Link:} \href{https://github.com/MadhuranS/capstone/pull/127}{link} \\
    \textbf{Summary:} Professor Mahalec wanted the ability to modify the contents of the streams in the diagram so that he would have better control of the system. He felt that the current implementation was not user driven enough and so by adding stream selection the user has more control. 
\subsection{Design and Design Documentation/Rev 0}

    \subsubsection*{Feedback Label: Add translation function that translates canvas data to required spec}
    \textbf{Source:} Professor Mahalec during Rev 0 work \\
    \textbf{Issue Link:} \href{https://github.com/MadhuranS/capstone/issues/139}{link} \\
    \textbf{PR Link:} \href{https://github.com/MadhuranS/capstone/pull/194}{link} \\
    \textbf{Summary:} Professor Mahalec wanted the data to be specific in a specific format so that his model could read the data in a certain way. To implement this we built a translator function that translates the data from canvas format to his format.  

    \subsubsection*{Feedback Label: Add team name and update description and extras}
    \textbf{Source:} Professor Smith for Rev 0 deadline \\
    \textbf{Issue Link:} \href{https://github.com/MadhuranS/capstone/issues/144}{link} \\
    \textbf{PR Link:} \href{https://github.com/MadhuranS/capstone/pull/171}{link} \\
    \textbf{Summary:} Professor Smith wanted us to provide a specific team name and update our project description and extras

    \subsubsection*{Feedback Label: Improve visibility for information in node ports}
    \textbf{Source:} Professor Mahalac for rev 0 \\
    \textbf{Issue Link:} \href{https://github.com/MadhuranS/capstone/issues/255}{link} \\
    \textbf{PR Link:} \href{https://github.com/MadhuranS/capstone/pull/295}{link} \\
    \textbf{Summary:} Professor Mahalac wanted better visibility into each ports info for each node and so we added tooltips and descriptive labels for the components that lacked info. 

    \subsubsection*{Feedback Label: Add table for material contents}
    \textbf{Source:} Professor Mahalac's Grad student Farbod for rev 0 \\
    \textbf{Issue Link:} \href{https://github.com/MadhuranS/capstone/issues/229}{link} \\
    \textbf{PR Link:} \href{https://github.com/MadhuranS/capstone/pull/263}{link} \\
    \textbf{Summary:} Farbod wanted an interface we could use to view stream contents and so we implemented an HTML table that was then revised into a react table to provide better data visibility. 

    \subsubsection*{Feedback Label: Add Docker setup to run the project}
    \textbf{Source:} Capstone TA Yiding for Rev 0\\
    \textbf{Issue Link:} \href{https://github.com/MadhuranS/capstone/issues/206}{link} \\
    \textbf{PR Link:} \href{https://github.com/MadhuranS/capstone/pull/210}{link} \\
    \textbf{Summary:} Our TA Yiding wanted a way to deploy the project in a simple manner and so we included a docker setup in our design. By dockerizing our app we made it easier for anyone to run the app regardless of technical experience. 

\subsection{VnV Plan and Report/Rev 1}

    \subsubsection*{Feedback Label: Write a code walkthrough so that the new team can get onboarded}
    \textbf{Source:} New team members who are taking on the project in May\\
    \textbf{Issue Link:} \href{https://github.com/MadhuranS/capstone/issues/393}{link} \\
    \textbf{PR Link:} \href{https://github.com/MadhuranS/capstone/pull/392}{link} \\
    \textbf{Summary:} New team members wanted a guide for the code to make it easier to understand and so we wrote out a code walkthrough. 

    \subsubsection*{Feedback Label: Remove out of scope requirement from docs}
    \textbf{Source:} Professor Mahalec during rev 1 docs meeting\\
    \textbf{Issue Link:} \href{https://github.com/MadhuranS/capstone/issues/394}{link} \\
    \textbf{PR Link:} \href{https://github.com/MadhuranS/capstone/pull/391}{link} \\
    \textbf{Summary:} Our authentication levels requirement was out of scope for our project and so Professor Mahalec indicated that we could remove this requirement from our docs. 

    \subsubsection*{Feedback Label: Make icon sizing data driven and dynamic}
    \textbf{Source:} Professor Mahalec Grad TA Raunak during rev 1 review meeting\\
    \textbf{Issue Link:} \href{https://github.com/MadhuranS/capstone/issues/348}{link} \\
    \textbf{PR Link:} \href{https://github.com/MadhuranS/capstone/pull/382}{link} \\
    \textbf{Summary:} Raunak instructed us that the sizez of icons on the diagram must be dynamic and user data driven so that the visuals of the diagram are also controlled by the user. We made this change as we felt that it helped improve useability. 

    \subsubsection*{Feedback Label: Add unit tests for backend endpoints}
    \textbf{Source:} TA Yiding after vnv plan meeting\\
    \textbf{Issue Link:} \href{https://github.com/MadhuranS/capstone/issues/395}{link} \\
    \textbf{PR Link:} \href{https://github.com/MadhuranS/capstone/pull/329}{link} \\
    \textbf{Summary:} As part of vnv plan meeting it was mentioned that for unit testing we could consider adding unit tests for our backend services since our frontend would be difficult to unit test directly. We complied by adding these unit tests to our endpoints. 
    
    \subsubsection*{Feedback Label: Add save confirmation check when exiting diagram}
    \textbf{Source:} Professor Mahalec during rev 1 meetings\\
    \textbf{Issue Link:} \href{https://github.com/MadhuranS/capstone/issues/292}{link} \\
    \textbf{PR Link:} \href{https://github.com/MadhuranS/capstone/pull/303}{link} \\
    \textbf{Summary:} Professor Mahalec wanted to add a confirmation check for saving the diagram when the user tries to exist the system. We agreed to make this change as we felt it made the system more robust and secure. 



\section{Challenge Level and Extras}

\subsection{Challenge Level}

Our project is a complex full-stack system, requiring backend services for version management and a sophisticated front-end for plant schema creation and modification. While not advanced due to limited novelty, it meets the criteria for a general challenge level due to the intricate development involved.

\subsection{Extras}

Since we are now handing over the project to a new set of developers at the end of the term, we decided that creating a code walkthrough and a user instructional video made more sense as extras. This is because our existing docs provide enough text based content for the user to follow and it would be easier for newer developers to learn how to use the system with a video. Similarly to help support newer developers we wrote out a code walkthrough so that they could learn the codebase faster. 

\section{Design Iteration (LO11 (PrototypeIterate))}

We started off our prototyping by building out the canvas component. We decided to start with this since its the core feature of our project and so we felt that it was important to start development with this feature and then work around it. After that, we met with out sponsor and his team to figure out what their requirements were for additonal features as well as for the development of the canvas as a whole. When our client indicated that they wanted the system to be data driven, we decided that we needed a user friendly way for users to input their system requirements and then we needed a method to translate those requirements into actual system data. This is how we devised the excel schemas which allowed users to easily enter data, we then built a python migration script that translated the excel into SQL data and then finally we built an API service that moved this sytem data to the client allowing the client to be data driven. 

Similarly, our users also felt that the current canvas implementation did not have enough interactiveness. Through our initial usability tests, they noticed there was certain data that was outside the scope of their control and certain components of the canvas that they lacked control over. To fix this we began a canvas overhaul that allowed us to add more interactive components to the canvas. This included the ability to freely modify data in each node and the ability to select a stream and modify its contents. It was clear at this point that a major theme for our users needs would be user driven features rather than system driven features. 

Moving on, after conducting further tests and meetings with these new data driven features, our sponsors noted that while they appreciated the fact that the system was user driven, they still wanted some constraints on what a user could change/not change so that the overall system was still robust and could not be broken by the users. This lead us to implementing checks on different modifications of data and also provude UI feedback to the user on the changes they made such that we control the amount of interactiveness the user could have with the client. By doing so we ensured that the overall system was flexible enough that the user had power but rigid enough so that the user could not break the system. 

Next we worked on adding final components that the user required. Once again going with the user driven interactiveness theme, these changes involved building a component in the canvas where the user could modify the systems material contents and time periods. These changes gave users the control they needed to have full control over the factory diagramming system to the point where it was almost like they were controlling a real factory which was the intent behind our project along. 

Finally, we reviewed the finished work with our user, conducted stress based usability tests to identify any bugs and then worked on cleaning up and fixing these bugs before eventually working on project handover. 


\section{Design Decisions (LO12)}

 \subsection{Limitations}
We did not have direct access to the ML model that would be responsible for running the simulations and producing the outputs of each factory plant since that model was also in development. So we had to use temporary proxy data in its place in order to mimic the ML model so that it was almost like we were using the ML model ourselves.

We also only had access to limited system data and so we had to ensure that the system data that we did have access to was fleshed out enough such that we could use it to test all of our project requirements. 

\subsection{Assumptions}
We had to assume that the user had some level of chemical engineering knowledge in order to work the system with minimal support. The system is complex so its impossible to provide enough visual guidance for the user if we dont make this assumption. Instead we assumed that all users have general understand of chemical engineering specifically related to factory plants and it's components.

We also had to assume that the user had the device specs to run the project on their browser. Its not project that requires extremely high specs but on very lowend machines theres a chance for lag on extremely scaled out diagrams. 

\subsection{Constraints}
Our budget was limited to purchasing just a reactflow premium license, everything else had to be free of charge. So we decided to use local implementations of the database rather than using external cloud hosted DB's, this was to avoid the potential charges that come with cloud hosted solutions. We also had to rely on free and open source libraries for any tooling in our system since we didn't have the budget to purchase any additional tooling.

The project timeframe was 8 months and so we had to pick and choose what we could develop so that we could meet the 8 month timeline.

Finally, our users are not software engineers and so a constraint was that our project had to be easy to use and deploy locally. So we had to dockerize our service to meet this constraint.  
 

\section{Economic Considerations (LO23)}


Our software was sponsored by Professor Mahalec and his grad team who were looking for a plant modelling software that could replace existing plant modelling software such as Aspen or Draw.io. The limitations for some of these other software is that Aspen has limits when it comes to scaling plant computations for multiple time periods while draw.io lacks the data driven features provided by our product. The product is intended to be sold by Professor Mahalec at a similar pricing point as Aspen (which costs \$20000 to \$50000 per year at a corporate license level although includes the computation model which is not part of our project) due to the similar but more expanded feature set compared to Aspen. Sales numbers would be determined by the number of corporations willing to pay for our time periods and data driven features. The cost to produce this product is free as its not dependant on any external licenses. Therefore our system is cost effective but still has the potential to bring in significant revenue due to the high demand  for powerful plant modelling systems. 


\section{Reflection on Project Management (LO24)}


\subsection{How Does Your Project Management Compare to Your Development Plan}


Yes we followed our development plan closely. We used discord and teams for communication and we relied heavily on our github project for project management. Our roles were also consistent however we did occasionally transfer over notetaking roles in order to avoid note taking burn out for different team members. We also followed our workflow plan closely by relying on a main branch and feature branch system that used PR's to move code into main. We did occasionally move away from our issue and pr templates since we felt that they were too rigid and on occassion we needed flexibility in these templates. Finally we did rely on CI/CD to produce our documents for the docs of our project. 

\subsection{What Went Well?}
Our use of the Github feature branch and main branch system went really well. By relying on feature branches for individual work we ensured that our work was well organized and scoped and that we didnt over complicate our individual commits. Similarly our PR process where each PR had to be reviewed by someone else made it so that all of our work was very visible to one another. Since we used squash commits and rebases in every PR, our commit history is also very linear and very clean which makes it easy for new developers to go back and review the commit history. 

\subsection{What Went Wrong?}

Our issue and PR templates may have been too rigid to cover every single use case for an issue or PR. For convience, we had to sometimes move away from our existing templates so that we could make clearer issues and pull request descriptions. Also sometimes relying on another reviewer would slow down work as in order for a change to be merged to main two people had to be available to merge it (one to write the PR and one to review it). 

\subsection{What Would you Do Differently Next Time?}
In the future, we would not limit ourselves to specific templates for PR's and issues so that there is some flexibility in project management. Also in the future we would have better traceability between our PR's and issues throughout the entire project. We only began adding that traceability toward the latter half of our project and so the first half it can be difficult to trace issues to PR's.

\section{Reflection on Capstone}

\subsection{Which Courses Were Relevant}

Courses related to design, documentation and architecture were all very relevant to this project. Design courses like 4HC3 helped teach good patterns and processes in UI/UX development. Courses like SFWRENG 3A04 helped us compose the overall project architecture. Finally courses like 3RA3 and 2AA4 gave us the documentation knowledge needed to properly document our work. 

\subsection{Knowledge/Skills Outside of Courses}


Alot of the implementation work was reliant on skills learned outside the course. For example how to do frontend programming with React and TypeScript was a skill we needed to learn ourselves. Similarly with developing the backend, we had to teach ourselves how to design REST API's. While the databases course did help with database design, we had to learn how to use the specific ORM we used and how to design schemas and queries using that ORM. 

\end{document}